% scaffold.tex -- ICLR 2027 draft. Numbers carry % source comments; prose follows the
% repository rules (no em dashes, colons, semicolons or lists in running prose).
\section{The Scaffold}
\label{sec-scaffold}

Narration-of-Thought is a system prompt that asks the model to reason through
the situation as a first-person narrative in five sections before it answers
(Figure~\ref{fig:prompt}). The first section names the decision-maker and what
they know. The second lists every person whose life the decision touches and
what is at stake for each. The third narrates, for each available action, its
consequences at least two steps into the future for each of those people. The
fourth states what remains uncertain about each projected future, and the fifth
commits to a decision and explains why that trajectory is preferable. The
control is standard chain-of-thought, a one-line instruction to think step by
step before answering. The prompt is identical on every model and task in this
paper, and it carries no task-specific examples.
% scripts/run_phase1_quartet.py PROMPTS["narrative_cot"], PROMPTS["standard_cot"]

The design follows the social account of objectivity. Organised criticism
across perspectives exposes errors that no single perspective corrects
\citep{longino1990science,mercier2011argumentative}, and a model answering a
person's account of their own conflict holds only that person's perspective.
The scaffold writes the missing perspectives into the trace before the answer
is given, and the same principle, placed outside a single trace, motivates the
collective of Section~\ref{sec-collective}, where the opposing perspectives are
separate agents with opposed interests.

\begin{figure}[t]
\centering
\fbox{\begin{minipage}{0.94\columnwidth}\small
\textbf{Narration-of-Thought (system prompt).}
You are a thoughtful advisor. When given an ethical dilemma, reason through it
as a five-part first-person narrative before giving your answer.\\[2pt]
Section 1 -- Protagonist: Name and briefly characterise the decision-maker (who
they are, their role, what they know).\\
Section 2 -- Stakeholders: List every person whose life intersects this
decision and state what is at stake for each.\\
Section 3 -- Consequences: For each available action, narrate its consequences
at least two steps into the future for each stakeholder.\\
Section 4 -- Uncertainty: State what remains genuinely uncertain about each
projected future.\\
Section 5 -- Decision: Commit to a specific decision and explain, within the
narrative frame, why that trajectory is preferable to the alternatives.\\[2pt]
Work through all five sections before giving your final answer.\\[4pt]
\textbf{Standard chain-of-thought (control).} You are a helpful assistant. Think
step by step, then give your answer.
\end{minipage}}
\caption{\textbf{The two system prompts compared throughout}, verbatim.}
\label{fig:prompt}
\end{figure}
